\section{Modélisation}
\label{sec:modelisation}

Cette section décrit le modèle en temps continu, avant toute discrétisation :
ce que représente chaque grandeur et ce que suppose chaque hypothèse. Le schéma
programmé fait l'objet de la section~\ref{sec:discretisation}. Les grandeurs
sont en unités SI, et les valeurs citées sont celles de la configuration de
référence \texttt{experiments/piscine.yaml}.

Le cadre est celui des modèles de foule à contraintes dures de Maury et
Venel~\cite{maury2008,maury2011} : chaque personne cherche à suivre une vitesse
souhaitée, et une contrainte d'exclusion corrige ce souhait dès que deux
personnes se touchent. Notre modèle s'en écarte sur un point, l'inertie. La
vitesse n'est pas égale à la vitesse souhaitée : elle la rejoint en un temps
fini, avec une accélération bornée par l'adhérence du sol. C'est par là
qu'entre l'effet du carrelage mouillé.

\subsection{État du système}
\label{sec:etat}

La foule compte $N$ personnes, représentées par des disques rigides dans le
plan. L'agent $i$ est décrit par sa position $\pos_i\in\mathbb{R}^2$
(\si{\metre}), sa vitesse $\vit_i\in\mathbb{R}^2$ (\si{\metre\per\second}) et
son rayon $r_i$ (\si{\metre}), fixé à \SI{0.2}{\metre} dans toutes nos
campagnes. On note $\pos=(\pos_1,\dots,\pos_N)\in\mathbb{R}^{2N}$ la
configuration de la foule.

Chaque agent porte aussi un état discret, qui indique où il en est de son
évacuation :
\[
\text{eau}\ \longrightarrow\ \text{franchissement}\ \longrightarrow\
\text{plage}\ \longrightarrow\ \text{sorti}.
\]
Seuls les agents dans l'eau ou sur la plage suivent la dynamique et les
contraintes décrites ci-dessous ; les deux autres états sont définis en
section~\ref{sec:franchissements}.

Le domaine est une emprise rectangulaire de \SI{48}{\metre} sur
\SI{20.5}{\metre}. Le bassin, de \SI{25}{\metre} sur \SI{12.5}{\metre}, est
bordé de plages de \SI{3}{\metre} à gauche, de \SI{4}{\metre} en haut et en bas,
et de \SI{20}{\metre} du côté de la porte. Cette porte unique, de
\SI{1.2}{\metre} de large, est percée au milieu du mur de droite. Le bord du
bassin est un mur pour les deux populations : on ne passe de l'eau à la plage
que par un site de franchissement. Les cent agents partent du bassin, soit une
densité initiale d'environ \num{0.32} personne par mètre carré.

\subsection{Vitesse souhaitée et champ de guidage}
\label{sec:guidage}

Chaque agent cherche à suivre la vitesse souhaitée
\begin{equation}
\label{eq:vitesse-souhaitee}
\vect{U}_i = s_i\,\vect{e}(\pos_i),
\end{equation}
de norme $s_i$ et de direction unitaire $\vect{e}$. La norme dépend du milieu
$m$ sous l'agent et du niveau de panique : $s_i=\kappa\,s_m$, où $s_m$ est la
vitesse de déplacement dans le milieu et $\kappa=1{,}5$ un facteur de panique.
Sur la plage, $s_i=1{,}5\times\SI{1.34}{\metre\per\second}\approx\SI{2.0}{\metre\per\second}$,
au-delà du seuil où Helbing et al.\ observent l'effet \emph{faster-is-slower}~\cite{helbing2000}.

\paragraph{Plus court temps de parcours.}
La direction suit le chemin le plus rapide vers la porte. Soit $\varphi(\pos)$
le temps minimal pour l'atteindre en se déplaçant à la vitesse locale
$F(\pos)=\kappa\,s_m$, la vitesse souhaitée du milieu sous $\pos$ ; il est
solution de l'équation eikonale
\begin{equation}
\label{eq:eikonale}
\norm{\nabla\varphi(\pos)}=\frac{1}{F(\pos)},
\qquad \varphi=0 \text{ sur la porte},
\qquad \vect{e}=-\frac{\nabla\varphi}{\norm{\nabla\varphi}} .
\end{equation}
Un obstacle est une région où $F=0$ : le temps de parcours y est infini, et le
champ le contourne. Comme $F$ dépend du milieu, $\varphi$ arbitre de lui-même
entre un chemin court mais lent et un chemin long mais rapide. L'équation est
résolue une fois pour toutes par Fast Marching~\cite{sethian1996}, sur une
grille de pas \SI{0.1}{\metre}.

L'eau et la plage ne communiquent que par les sites de franchissement, d'où
deux champs. Sur la plage, $\varphi_{\text{plage}}$ est le temps jusqu'à la
porte, le bassin étant un obstacle. Dans l'eau, chaque site $k$ devient une
source dont le temps initial vaut
\begin{equation}
\label{eq:phi-eau}
\varphi_{\text{eau}}(\vect{a}_k)=T_k+\varphi_{\text{plage}}(\vect{b}_k),
\end{equation}
où $\vect{a}_k$ est le point d'attente du site dans l'eau, $\vect{b}_k$ son
point d'arrivée sur la plage et $T_k$ la durée du franchissement.
$\varphi_{\text{eau}}$ est donc le temps \emph{total} jusqu'à la porte, et chaque
point de l'eau choisit le site qui le minimise : nager plus longtemps vers une
échelle rapide, ou se hisser sur le rebord tout proche mais lent. Le partage du
flux entre rebord et échelles n'est pas imposé ; il émerge des vitesses, des
durées et de la géométrie. C'est la partie « extraction répartie » de la
problématique. Le champ est statique : il ignore l'attente devant un site occupé
ou devant une porte encombrée. C'est une limite assumée du modèle.

\paragraph{Régulation vitesse--espacement.}
Derrière un autre agent, la norme souhaitée est réduite par la fonction de
vitesse optimale de Tordeux et al.~\cite{tordeux2016} :
\begin{equation}
\label{eq:espacement}
s_i \longleftarrow \min\bigl(s_i,\ \max(0,\ g_i/T)\bigr),
\end{equation}
où $g_i$ est le jeu entre le disque de $i$ et celui de son meneur, et
$T=\SI{0.35}{\second}$ un temps d'espacement. Le meneur est le plus proche des
agents $j$ situés devant $i$, dans un couloir de demi-largeur $r_i+r_j$ autour
de $\vect{e}(\pos_i)$, et plus avancés vers la sortie : $\varphi(\pos_j)<\varphi(\pos_i)$.
Dans la piscine, il est cherché du même côté du bord du bassin. La direction
n'est pas modifiée, et l'exclusion reste imposée par la seule contrainte de la
section~\ref{sec:exclusion}. Le choix de cette règle et l'étalonnage de $T$ sont
présentés en section~\ref{sec:e0}.

\subsection{Les trois milieux}
\label{sec:milieux}

Le sol agit sur la façon dont l'agent rejoint sa vitesse souhaitée. Hors
contact, la dynamique d'un agent s'écrit
\begin{equation}
\label{eq:relaxation}
\frac{\mathrm{d}\pos_i}{\mathrm{d}t}=\vit_i,
\qquad
\frac{\mathrm{d}\vit_i}{\mathrm{d}t}
=\Pi_{a_{\max}(\pos_i)}\!\left(\frac{\vect{U}_i-\vit_i}{\tau(\pos_i)}\right),
\qquad
\Pi_{a}(\vect{w})=\vect{w}\,\min\Bigl(1,\frac{a}{\norm{\vect{w}}}\Bigr).
\end{equation}
Le temps de relaxation $\tau$ fixe la rapidité de la correction : sans
saturation, l'écart entre $\vit_i$ et $\vect{U}_i$ décroît comme
$\mathrm{e}^{-t/\tau}$. L'accélération maximale $a_{\max}$ borne la traction que
le sol peut fournir ; demander davantage, c'est glisser. La borne porte sur la
norme et non sur chaque composante : un agent qui freine en diagonale freine
ainsi dans la direction voulue, seulement moins fort.

\paragraph{Adhérence.}
La force horizontale que le sol transmet au pied est bornée par le frottement :
$\norm{M\vect{a}}\leq\mu M g$ pour une personne de masse $M$, soit
$a_{\max}=\mu g$, où $\mu$ est le coefficient de frottement disponible. Sur des
sols de piscine mouillés, parcourus pieds nus par 70 sujets,
Kendzior~\cite{kendzior2023} trouve qu'un coefficient dynamique inférieur ou
égal à \num{0.34} est perçu comme glissant, et un coefficient supérieur ou égal
à \num{0.60} comme non glissant. Nous retenons ces deux seuils comme coefficients
disponibles du carrelage mouillé et du sol sec, d'où
$a_{\max}=\SI{3.3}{\metre\per\second\squared}$ et
\SI{5.9}{\metre\per\second\squared}. C'est une hypothèse : un seuil de
perception n'est pas une mesure du frottement sous une foule qui court. Dans
l'eau, la poussée d'Archimède réduit le poids apparent, d'environ moitié dans
un petit bain, d'où $0{,}34\,g\times 0{,}5\approx\SI{1.67}{\metre\per\second\squared}$.

\paragraph{Vitesses et relaxation.}
Dans un petit bain de \SI{0.7}{\metre} d'eau calme, la vitesse de progression
suit $s=0{,}53\,A^{-0{,}19}$ avec $A=d^2/2$ (\si{\metre\squared}), où $d$ est la profondeur d'eau, d'après
Shirvani et al.~\cite{shirvani2020}, soit \SI{0.69}{\metre\per\second}. Sur la
plage, la vitesse de marche libre est celle de
Weidmann~\cite{weidmann1993}, \SI{1.34}{\metre\per\second}, sur sol mouillé
comme sur sol sec. Le temps de relaxation vaut \SI{0.5}{\second} sur les trois
sols, comme chez Helbing et al.~\cite{helbing2000} : aucune source ne donne un
$\tau$ propre au carrelage mouillé, et un $\tau$ commun fait que le mouillé ne
diffère du sec \emph{que} par l'adhérence, ce qui isole l'effet étudié. Le
tableau~\ref{tab:milieux} résume ces valeurs.

\begin{table}[ht]
\centering
\caption{Paramètres des trois milieux, avant application du facteur de panique
$\kappa=1{,}5$ à la vitesse. Le carrelage mouillé est une bande de
\SI{3}{\metre} autour du bassin, remplacée par du sol sec dans la campagne
témoin E1.}
\label{tab:milieux}
\begin{tabular}{lrrrl}
\toprule
Milieu & $s_m$ (\si{\metre\per\second}) & $\tau$ (\si{\second})
& $a_{\max}$ (\si{\metre\per\second\squared}) & Sources \\
\midrule
Eau, petit bain & 0,69 & 0,5 & 1,67 & \cite{shirvani2020}, hypothèse de poids apparent \\
Carrelage mouillé & 1,34 & 0,5 & 3,3 & \cite{weidmann1993,kendzior2023} \\
Sol sec & 1,34 & 0,5 & 5,9 & \cite{weidmann1993,kendzior2023} \\
\bottomrule
\end{tabular}
\end{table}

\paragraph{Où se loge l'inertie.}
À la vitesse souhaitée de panique, $U\approx\SI{2.01}{\metre\per\second}$,
repartir de l'arrêt demande $U/\tau\approx\SI{4.0}{\metre\per\second\squared}$.
C'est moins que les \SI{5.9}{\metre\per\second\squared} du sol sec, mais plus que
les \SI{3.3}{\metre\per\second\squared} du mouillé : seul le mouillé écrête
l'accélération, et un agent n'y peut ni freiner ni tourner aussi vite qu'il le
voudrait. Ce contraste porte l'effet d'inertie de la problématique.

\subsection{Principe d'exclusion}
\label{sec:exclusion}

Deux personnes ne peuvent pas se recouvrir, et chacune reste dans la partie du
domaine où elle se trouve. Les configurations admissibles forment l'ensemble
\begin{equation}
\label{eq:admissible}
\admissible=\bigl\{\pos\in\mathbb{R}^{2N} :
D_{ij}(\pos)=\norm{\pos_j-\pos_i}-r_i-r_j\geq 0\ \ \forall\, i<j,\
\pos_i\in\Omega_i\bigr\},
\end{equation}
où $\Omega_i$ est la région accessible au centre de l'agent $i$ : le bassin
réduit de son rayon pour un agent dans l'eau ; pour un agent sur la plage,
l'ensemble des points de la plage et du couloir de la porte situés à une
distance au moins $r_i$ de tout mur et du bord du bassin. Près d'un coin, de
mur ou de bassin, ce bord est donc un arc de cercle de rayon $r_i$. Une
première version du code remplaçait cet arc par un angle vif aux montants de
la porte : la foule s'y bloquait en arches statiques, sans que ce soit un
comportement de foule.

\paragraph{$\admissible$ n'est pas convexe.}
C'est le point délicat du modèle. Prenons deux disques de rayon $r$ tangents,
centrés en $(-r,0)$ et $(r,0)$ : la configuration est admissible, et celle où
les deux disques échangent leurs places l'est aussi. Leur milieu place pourtant
les deux centres à l'origine, avec $D_{12}=-2r<0$. Chaque contrainte
$D_{ij}\geq0$ décrit l'extérieur d'une boule de $\mathbb{R}^{2N}$, et le
domaine de la plage, un anneau autour du bassin, n'est pas convexe non plus. La
projection sur $\admissible$ n'est donc en général ni unique ni calculable
directement ; la section~\ref{sec:discretisation} explique comment on la
remplace par une suite de projections sur des ensembles convexes.

\paragraph{Contact et choc mou.}
Dans le modèle de Maury et Venel, du premier ordre, la vitesse effective est la
projection de la vitesse souhaitée sur le cône des vitesses admissibles :
\[
\frac{\mathrm{d}\pos}{\mathrm{d}t}=P_{C_{\pos}}(\vect{U}),
\qquad
C_{\pos}=\bigl\{\vect{w}\in\mathbb{R}^{2N} :
D_{ij}(\pos)=0\ \Rightarrow\ \nabla D_{ij}(\pos)\cdot\vect{w}\geq 0\bigr\}.
\]
Notre modèle garde l'inertie de l'équation~\eqref{eq:relaxation}. Au contact,
la vitesse est corrigée de la même façon : la composante qui ferait se
recouvrir deux disques est annulée, l'autre est conservée. Le choc est donc
\emph{mou}, sans rebond, et sans frottement entre agents puisque la correction
agit le long de la ligne des centres. Le modèle ne contient aucune force de
contact : seule la géométrie limite la densité. Pour des disques de rayon
\SI{0.2}{\metre}, l'empilement hexagonal compact atteint
$1/(2\sqrt{3}\,r^2)\approx\num{7.2}$ personnes par mètre carré, au-dessus de la
densité maximale de \num{5.4} personnes par mètre carré retenue par
Weidmann~\cite{weidmann1993}.

\subsection{Franchissements à capacité limitée}
\label{sec:franchissements}

Le bord du bassin est découpé en \emph{sites de franchissement}. Le rebord
compte un site tous les \SI{0.6}{\metre}, soit 118 sites, et deux échelles de
\SI{0.5}{\metre} de large sont placées en haut et en bas, aux neuf dixièmes de
la longueur du bassin, du côté de la porte. Chaque site $k$ a un point d'attente $\vect{a}_k$ dans
l'eau, un point d'arrivée $\vect{b}_k$ sur la plage, une durée $T_k$ et une
capacité d'un agent :
\begin{itemize}
\item un agent dans l'eau à moins de \SI{0.25}{\metre} du point d'attente d'un
  site libre l'occupe. Quand plusieurs couples (agent, site) sont possibles, ils
  sont servis du plus proche au plus éloigné, si bien que l'attribution dépend
  des positions et non de la numérotation des agents ;
\item l'agent est immobilisé sur le bord pendant $T_k$, hors de la contrainte
  d'exclusion : \SI{4.0}{\second} pour se hisser sur le rebord,
  \SI{2.0}{\second} pour monter une échelle ;
\item il est ensuite réintroduit sur la plage en $\vect{b}_k$ dès que ce point
  est libre de tout agent actif. Le site reste occupé jusque-là.
\end{itemize}
Les durées ne reposent sur aucune source : ce sont des hypothèses, dont
l'influence doit être mesurée par balayage (E4). Un agent dont le centre franchit la ligne de la porte
passe dans l'état \emph{sorti}, qui est absorbant : il quitte la contrainte
d'exclusion et le guidage, et ne peut plus gêner la porte. Son instant de
sortie est interpolé sur la ligne. Le nombre d'agents dans chacun des quatre
états change, mais leur somme reste égale à $N$ ; cette conservation est
vérifiée à chaque pas (section~\ref{sec:validation}).

Un module de chutes (glissade, propagation en domino, relevage et contournement
d'une personne à terre) existe dans le code. Il est désactivé dans toutes les
simulations de ce rapport et ses taux ne sont pas calibrés : il n'intervient
dans aucun résultat présenté.

\subsection{Hypothèses principales}

Les hypothèses à garder en tête pour lire les résultats sont les suivantes :
\begin{itemize}
\item le modèle est plan, et chaque personne est un disque rigide de même
  rayon ;
\item la population est homogène : mêmes vitesses souhaitées, même
  relaxation, même facteur de panique pour tous ;
\item le milieu est constant par morceaux et lu sous le centre de l'agent ;
\item le guidage est statique et ne voit ni les autres agents ni l'attente aux
  sites ;
\item les contacts sont des chocs mous sans frottement, et aucune force de
  compression n'est modélisée ;
\item les durées de franchissement et le poids apparent dans l'eau sont des
  hypothèses non sourcées, dont l'influence reste à mesurer par balayage ;
\item le temps d'espacement $T$ est étalonné sur un goulot sec, à la marche
  normale ; il est appliqué tel quel au facteur de panique $1{,}5$ et dans
  l'eau. Ce transfert n'est pas neutre : au facteur $1{,}5$, le même goulot
  écoule \num{2.04} au lieu de \num{1.84}~personnes par mètre et par seconde
  (section~\ref{sec:e1}).
\end{itemize}
